\section{Conclusion}
\label{sec:conclusion}

This project set out to answer how a software system can automate medical consultation documentation using local AI assistance while preserving privacy, security, and architectural flexibility. Med Flow demonstrates one practical answer: sensitive speech-to-text and LLM processing can run locally, AI output can remain editable until doctor approval, and the documentation workflow can be structured as replaceable components rather than as one tightly coupled application.

The implemented system supports the core outpatient workflow from appointment context to patient follow-up. A doctor can start a consultation, record audio, receive a transcript, generate a structured clinical draft, inspect Suggestive Mode alerts, edit the report, approve or reject the result, export a PDF, and send approved follow-up communication. The patient and administrator roles support this workflow through patient-facing access, booking context, configuration, and auditability, while the doctor remains responsible for clinical sign-off.

The architecture is intentionally component-based. FastAPI routes, application services, domain contracts, infrastructure adapters, AI services, persistence layers, and notification tools are separated so the system can evolve without forcing every change through a monolithic structure. The dual-database design supports both reliable relational records and flexible document-oriented AI artifacts, while the replaceable boundaries around transcription, LLM generation, scheduling, identity, PDF generation, and email delivery keep the system open to future integration.

The validation package supports the core engineering claims of the project. Automated tests, stress tests, benchmark tests, hallucination checks, ASR artifact checks, manual UI validation, and hardware-envelope testing show that the implementation is testable, measurable, and suitable as a foundation for further development. \tabref{tab:requirements-coverage} shows that the central must-have requirements were implemented, while deployment security, retention policy, cross-doctor sharing, production identity integration, and broader end-to-end automation remain clear future-work items.

The main limitation is that Med Flow is still a semester-scale prototype rather than a clinically deployed product. It has not been validated with real patients, real clinical audio corpora, production security controls, or integration with national identity and booking infrastructure. Those limitations do not weaken the core result; instead, they define the next engineering steps needed before a system like this could move from a local demonstrator toward healthcare deployment.

\begin{summarybox}
Med Flow is best understood as a controlled AI documentation workflow: local AI accelerates transcription, drafting, and review, modular architecture keeps the system adaptable, validation provides evidence for the implemented behavior, and the doctor remains the final decision-maker for clinical approval, PDF export, and patient communication.
\end{summarybox}